% =====================================================================
% ANNEXES
% =====================================================================

% =====================================================================
\section{Annexe A~--- Fiches des cas d'utilisation}
\label{annexe:ucfiches}
% =====================================================================

Cette annexe complète la section~\ref{subsec:uc-fiches} du chapitre~5 en fournissant les fiches descriptives des 24 cas d'utilisation non développés dans le corps du document.

\subsection*{Phase 1~--- Project Management}

\ucfiche{UC-01}{Créer un projet de migration}
  {Administrator}
  {Utilisateur authentifié avec le rôle Administrator}
  {L'administrateur souhaite initier un nouveau projet de migration de données}
  {\textbf{1.} L'administrateur accède au tableau de bord global. \textbf{2.} Il clique sur ``Nouveau projet''. \textbf{3.} Il renseigne les informations~: nom, Industry, Target platform, Source system. \textbf{4.} Il valide la création. \textbf{5.} La plateforme initialise les huit phases et génère l'identifiant projet.}
  {Nom de projet déjà existant~: message d'erreur et suggestion d'un nom disponible. Champs obligatoires manquants~: validation inline.}
  {Projet créé avec les huit phases initialisées~; l'administrateur est automatiquement ajouté comme membre owner}

\ucfiche{UC-02}{Configurer les paramètres du projet}
  {Administrator}
  {Projet créé (UC-01)~; au moins un membre de l'équipe défini (UC-03)}
  {L'administrateur souhaite modifier ou affiner les paramètres d'un projet existant}
  {\textbf{1.} L'administrateur ouvre les paramètres du projet. \textbf{2.} Il modifie les informations disponibles~: description, dates cibles, options d'affichage. \textbf{3.} Il sauvegarde les modifications. \textbf{4.} Un historique des changements est enregistré.}
  {Champs protégés en mode lecture seule si le projet est verrouillé~: notification à l'administrateur. Sauvegarde concurrente~: résolution par timestamp.}
  {Paramètres mis à jour~; historique des modifications disponible dans les logs projet}

\ucfiche{UC-04}{Suivre l'avancement du projet}
  {Migration Designer, Cutover Manager}
  {Projet actif avec au moins une phase démarrée}
  {L'utilisateur consulte le tableau de bord pour suivre l'avancement global du projet}
  {\textbf{1.} L'utilisateur ouvre le tableau de bord du projet. \textbf{2.} Il visualise le statut de chaque phase (Non démarré, En cours, Terminé, Bloqué). \textbf{3.} Il consulte les métriques d'avancement~: pourcentage de complétion, jalons atteints, alertes en cours. \textbf{4.} Il peut naviguer directement vers une phase en cliquant dessus.}
  {Données non disponibles pour une phase non initialisée~: indicateur grisé. Alertes critiques~: bandeau rouge en haut du tableau de bord.}
  {Vue d'avancement consultée~; aucune modification de l'état du projet}

\ucfiche{UC-05}{Archiver un projet finalisé}
  {Administrator}
  {Projet en statut ``Completed''~; validation go-live effectuée (UC-29)}
  {L'administrateur archive le projet pour libérer les ressources actives et conserver le dossier}
  {\textbf{1.} L'administrateur sélectionne le projet finalisé. \textbf{2.} Il déclenche l'action ``Archiver''. \textbf{3.} La plateforme génère une archive complète~: documents, bundles, rapports, logs. \textbf{4.} Le projet passe en mode lecture seule~; les ressources computationnelles sont libérées.}
  {Projet avec des tâches cutover ouvertes~: archivage refusé jusqu'à clôture complète. Erreur d'archivage~: rollback et notification.}
  {Projet archivé en lecture seule~; archive téléchargeable pendant 12 mois}

\subsection*{Phase 2~--- Data Scope}

\ucfiche{UC-06}{Définir le périmètre des données}
  {Data Analyst, Migration Designer}
  {Projet initialisé (UC-01)~; phase Data Scope débloquée}
  {Le Data Analyst définit les entités sources, intermédiaires et cibles du périmètre de migration}
  {\textbf{1.} Le Data Analyst accède à la phase Data Scope. \textbf{2.} Il importe ou crée manuellement les définitions SDS (entités sources), TDS (entités cibles) et IDS (entités intermédiaires). \textbf{3.} Pour chaque entité, il renseigne les attributs~: nom, colonnes, types, clés. \textbf{4.} Il sauvegarde~; une version du data scope est créée.}
  {Entité en doublon~: avertissement et possibilité de fusion. Format invalide~: message d'erreur détaillé avec localisation.}
  {Data scope versionnée~; entités disponibles pour les phases suivantes~; verrouillage activable (UC-07)}

\ucfiche{UC-07}{Verrouiller une section du data scope}
  {Data Analyst}
  {Data scope défini (UC-06)~; section non verrouillée par un autre utilisateur}
  {Le Data Analyst verrouille une section du data scope pour l'éditer exclusivement}
  {\textbf{1.} Le Data Analyst sélectionne une entité ou un groupe d'entités. \textbf{2.} Il clique sur ``Checkout''. \textbf{3.} La plateforme vérifie qu'aucun verrou concurrent n'existe. \textbf{4.} Le verrou est posé~; les autres utilisateurs voient la section en lecture seule. \textbf{5.} Après édition, le Data Analyst effectue le ``Check-in'' pour libérer le verrou.}
  {Section déjà verrouillée~: affichage du nom du détenteur et heure de prise du verrou. Verrou expiré (> 4h)~: relâchement automatique après confirmation.}
  {Verrou libéré~; nouvelle version du data scope enregistrée~; section disponible pour d'autres éditeurs}

\ucfiche{UC-08}{Exporter le data scope}
  {Data Analyst}
  {Data scope finalisé et validé (UC-06)}
  {Le Data Analyst exporte le périmètre vers le format natif ADMP pour utilisation dans les outils de conception}
  {\textbf{1.} Le Data Analyst accède à la phase Data Scope. \textbf{2.} Il sélectionne le format d'export~: ADMP Native (Excel/Word) ou JSON structuré. \textbf{3.} La plateforme génère les fichiers correspondants (SDS, TDS, IDS). \textbf{4.} L'export est téléchargeable directement ou envoyé par email.}
  {Entités incomplètes~: avertissement avec liste des champs manquants avant export. Timeout d'export~: retry automatique.}
  {Fichiers SDS/TDS/IDS disponibles en téléchargement~; date d'export enregistrée dans les métadonnées}

\subsection*{Phase 3~--- Extract}

\ucfiche{UC-09}{Configurer une source d'extraction}
  {Data Analyst}
  {Phase Extract initialisée~; data scope défini (UC-06)}
  {Le Data Analyst configure les paramètres de connexion à une source de données à extraire}
  {\textbf{1.} Le Data Analyst accède à la phase Extract. \textbf{2.} Il ajoute une nouvelle source~: type (Oracle, SQL Server, CSV, VSAM), paramètres de connexion (hôte, port, base). \textbf{3.} Il teste la connexion. \textbf{4.} Il associe la source aux entités correspondantes du data scope. \textbf{5.} Il sauvegarde la configuration.}
  {Test de connexion échoué~: message d'erreur technique et suggestions de diagnostic. Entité non trouvée dans la source~: avertissement et option d'alignement manuel.}
  {Source configurée et associée aux entités~; prête pour l'import de fichiers (UC-10)}

\subsection*{Phase 4~--- Profiling \& Cleansing}

\ucfiche{UC-12}{Générer un rapport de profilage}
  {Data Analyst}
  {Audit de données lancé (UC-11)~; résultats disponibles}
  {Le Data Analyst génère un rapport formel de profilage statistique par entité}
  {\textbf{1.} Le Data Analyst accède aux résultats d'audit (UC-11). \textbf{2.} Il sélectionne les entités à inclure dans le rapport. \textbf{3.} Il choisit le format~: PDF, Excel, HTML interactif. \textbf{4.} La plateforme génère le rapport avec~: distributions, taux de nullité, cardinalités, exemples d'anomalies. \textbf{5.} Le rapport est archivé dans le dossier projet.}
  {Données insuffisantes pour une entité~: section correspondante marquée ``Données insuffisantes''. Génération longue~: notification par email.}
  {Rapport archivé~; disponible pour partage avec le Business Reviewer~; base pour les règles de nettoyage (UC-13)}

\ucfiche{UC-13}{Documenter les règles de nettoyage}
  {Data Analyst}
  {Rapport de profilage disponible (UC-12)~; anomalies identifiées}
  {Le Data Analyst documente les règles de nettoyage à appliquer aux données identifiées comme non conformes}
  {\textbf{1.} Le Data Analyst consulte les anomalies du rapport de profilage. \textbf{2.} Pour chaque anomalie, il crée une règle de nettoyage~: type (substitution, suppression, normalisation), expression, entités concernées. \textbf{3.} Il documente la justification métier. \textbf{4.} Il sauvegarde les règles~; elles sont versionnées et traçables.}
  {Règle en conflit avec une règle existante~: avertissement et possibilité de remplacement ou fusion. Expression invalide~: validation en temps réel.}
  {Règles de nettoyage documentées~; prêtes pour soumission (UC-15) et validation métier (UC-16)}

\ucfiche{UC-14}{Valider la conformité qualité des données}
  {Business Reviewer}
  {Règles de nettoyage appliquées~; rapport de profilage post-nettoyage disponible}
  {Le Business Reviewer évalue la conformité des données nettoyées par rapport aux exigences métier}
  {\textbf{1.} Le Business Reviewer reçoit une notification de disponibilité. \textbf{2.} Il consulte le rapport de profilage post-nettoyage. \textbf{3.} Il vérifie la conformité entité par entité. \textbf{4.} Pour chaque entité, il renseigne son avis (Conforme, Non conforme + commentaire). \textbf{5.} Il soumet sa validation.}
  {Entité non conforme~: commentaire obligatoire~; renvoi au Data Analyst pour correction. Délai dépassé (> 48h sans réponse)~: escalade automatique vers le Migration Designer.}
  {Validation métier enregistrée~; phase Profiling débloquée pour passage à Transform si toutes entités conformes}

\ucfiche{UC-15}{Soumettre les règles pour validation}
  {Data Analyst}
  {Règles de nettoyage documentées (UC-13)~; phase Profiling en cours}
  {Le Data Analyst soumet les règles de nettoyage au Business Reviewer pour approbation}
  {\textbf{1.} Le Data Analyst finalise les règles de nettoyage dans l'éditeur. \textbf{2.} Il clique sur ``Soumettre pour validation''. \textbf{3.} La plateforme notifie le ou les Business Reviewers assignés. \textbf{4.} Les règles passent en statut ``En attente de validation''.}
  {Aucun Business Reviewer assigné~: blocage avec message demandant l'assignation via UC-03. Règles incomplètes~: validation inline avant soumission.}
  {Règles soumises~; Business Reviewer notifié~; Data Analyst informé de la prise en charge}

\ucfiche{UC-16}{Approuver les règles de nettoyage}
  {Business Reviewer}
  {Règles soumises (UC-15)~; Business Reviewer notifié}
  {Le Business Reviewer approuve ou rejette les règles de nettoyage soumises par le Data Analyst}
  {\textbf{1.} Le Business Reviewer ouvre la liste des règles en attente. \textbf{2.} Il examine chaque règle~: description, entités concernées, justification métier. \textbf{3.} Il approuve (vert) ou rejette (rouge + commentaire) chaque règle. \textbf{4.} Il soumet sa décision. \textbf{5.} Le Data Analyst est notifié.}
  {Règle ambiguë~: possibilité de demander une clarification au Data Analyst via commentaire. Approbation partielle~: les règles approuvées sont validées~; les rejetées retournent en édition.}
  {Règles approuvées enregistrées~; règles rejetées renvoyées au Data Analyst~; journal d'approbation tracé}

\subsection*{Phase 5~--- Transform}

\ucfiche{UC-19}{Configurer les tables de transcodification}
  {Migration Designer}
  {Data scope finalisé~; phase Transform débloquée}
  {Le Migration Designer configure les tables de transcodification (TT) pour une unité de migration}
  {\textbf{1.} Le Migration Designer accède à la gestion des transco. \textbf{2.} Il importe un fichier TT (XML) ou crée les entrées manuellement. \textbf{3.} Il définit les clés et valeurs de mapping. \textbf{4.} Il associe la TT aux unités de migration qui l'utilisent. \textbf{5.} Il valide et sauvegarde.}
  {Clé dupliquée dans la TT~: avertissement et proposition de merge. Format XML invalide~: erreur de validation avant import.}
  {Tables de transco disponibles pour les règles RC/DP~; vues broadcast générées automatiquement lors de la compilation}

\ucfiche{UC-20}{Valider le bundle généré}
  {Test Engineer}
  {Bundle généré et disponible (UC-18)}
  {Le Test Engineer exécute des tests de validation sur le bundle généré avant déploiement en production}
  {\textbf{1.} Le Test Engineer accède au bundle dans la phase Transform. \textbf{2.} Il lance l'exécution sur un cluster de test Databricks. \textbf{3.} Il vérifie les compteurs (lignes lues, écrites, rejetées). \textbf{4.} Il consulte les rapports d'intégrité Excel générés. \textbf{5.} Il documente le résultat~: approuvé ou rejeté avec détail des anomalies.}
  {Erreur d'exécution sur le cluster~: log Spark UI disponible pour diagnostic. Compteurs inattendus~: comparaison avec la baseline de référence.}
  {Résultat de validation enregistré~; si approuvé, bundle prêt pour déploiement (UC-21)}

\ucfiche{UC-21}{Publier le bundle sur Databricks}
  {Migration Designer}
  {Bundle validé (UC-20)~; cluster Databricks de production accessible}
  {Le Migration Designer déploie le bundle validé sur le cluster Databricks de production}
  {\textbf{1.} Le Migration Designer sélectionne le bundle validé. \textbf{2.} Il choisit le cluster cible et les paramètres de déploiement. \textbf{3.} La plateforme appelle l'API Databricks~: upload du .whl, installation de la librairie. \textbf{4.} Un test de smoke est exécuté automatiquement (import du module). \textbf{5.} La plateforme confirme le déploiement réussi.}
  {Cluster indisponible~: notification et retry programmé. Test de smoke échoué~: rollback de l'installation et notification.}
  {Bundle installé sur le cluster~; disponible pour l'exécution du chargement (UC-23)}

\subsection*{Phase 6~--- Load}

\ucfiche{UC-22}{Configurer le chargement cible}
  {Migration Designer}
  {Bundle déployé (UC-21)~; système cible accessible (Delta Lake / Unity Catalog)}
  {Le Migration Designer configure les paramètres d'exécution du chargement vers le système cible}
  {\textbf{1.} Le Migration Designer accède à la phase Load. \textbf{2.} Il sélectionne le bundle déployé. \textbf{3.} Il renseigne les paramètres run-time~: catalog, schéma, paramètres métier, mode d'exécution (full/delta). \textbf{4.} Il configure l'ordre d'exécution des unités. \textbf{5.} Il sauvegarde la configuration de run.}
  {Catalog ou schéma introuvable~: vérification via Unity Catalog API~; message d'erreur détaillé. Paramètre manquant~: validation inline.}
  {Configuration de chargement sauvegardée~; prête pour exécution (UC-23)}

\subsection*{Phase 7~--- Validation}

\ucfiche{UC-24}{Définir les cas de test}
  {Test Engineer}
  {Chargement exécuté (UC-23)~; données présentes dans le système cible}
  {Le Test Engineer définit les cas de test et les critères d'acceptation pour la réconciliation}
  {\textbf{1.} Le Test Engineer accède à la phase Validation. \textbf{2.} Il crée les cas de test~: entités à vérifier, type de contrôle (count, spot-check, intégrité), valeurs attendues. \textbf{3.} Il associe chaque cas à une entité du data scope. \textbf{4.} Il sauvegarde et documente la justification des critères.}
  {Entité non chargée~: avertissement~; cas de test créé en statut ``En attente de données''. Critère sans valeur attendue~: validation requise avant exécution.}
  {Cas de test définis~; prêts pour exécution de la réconciliation (UC-25)}

\ucfiche{UC-25}{Exécuter les contrôles de réconciliation}
  {Test Engineer}
  {Cas de test définis (UC-24)~; données sources et cibles disponibles}
  {Le Test Engineer exécute les contrôles de réconciliation pour comparer les données source et cible}
  {\textbf{1.} Le Test Engineer sélectionne les cas de test à exécuter. \textbf{2.} Il lance la réconciliation. \textbf{3.} La plateforme interroge les systèmes source et cible~: comptages, spot-checks, vérifications d'intégrité. \textbf{4.} Les résultats sont comparés aux valeurs attendues. \textbf{5.} Les écarts sont classifiés (bloquant / majeur / mineur).}
  {Connexion source perdue~: résultats partiels marqués~; retry programmé. Timeout de contrôle~: reprise incrémentale possible.}
  {Résultats disponibles~; écarts classifiés~; rapport de réconciliation générable (UC-26)}

\subsection*{Phase 8~--- Cutover}

\ucfiche{UC-27}{Créer le plan de cutover}
  {Cutover Manager}
  {Rapport de réconciliation approuvé~; phase Cutover débloquée}
  {Le Cutover Manager crée un plan de cutover structuré pour orchestrer la bascule en production}
  {\textbf{1.} Le Cutover Manager accède à la phase Cutover. \textbf{2.} Il crée un plan~: nom, description, date cible de go-live. \textbf{3.} Il ajoute les tâches~: titre, responsable, durée estimée, ordre, dépendances. \textbf{4.} Il définit les jalons et critères go/no-go. \textbf{5.} Il soumet le plan pour revue par l'équipe.}
  {Dépendance circulaire entre tâches~: détection et avertissement~; résolution requise. Date cible incompatible avec les dépendances~: recalcul de la date proposé.}
  {Plan de cutover créé et versionné~; partagé avec l'équipe~; tâches assignées aux responsables}

\ucfiche{UC-28}{Exécuter les tâches de cutover}
  {Cutover Manager}
  {Plan de cutover approuvé (UC-27)~; équipe notifiée~; fenêtre de cutover ouverte}
  {Le Cutover Manager pilote l'exécution des tâches du plan de cutover en temps réel}
  {\textbf{1.} Le Cutover Manager ouvre la vue d'exécution du plan. \textbf{2.} Il démarre les tâches selon leur séquence et dépendances. \textbf{3.} Chaque responsable marque ses tâches comme terminées au fil de l'exécution. \textbf{4.} Le Cutover Manager suit l'avancement global sur le tableau de bord. \textbf{5.} Les incidents sont documentés en temps réel.}
  {Tâche bloquée ou en retard~: alerte immédiate avec escalade optionnelle. Incident critique~: activation du plan de rollback.}
  {Toutes les tâches complétées~; journal d'exécution horodaté~; bascule technique effectuée~; prêt pour go-live (UC-29)}

\subsection*{Modules transverses}

\ucfiche{UC-30}{Gérer la base de connaissances}
  {Administrator}
  {Accès Administrator à la plateforme}
  {L'administrateur met à jour la base de connaissances utilisée par les agents IA et les équipes}
  {\textbf{1.} L'administrateur accède au module Apprentissage. \textbf{2.} Il crée, modifie ou supprime des articles~: titre, contenu (Markdown), tags, catégorie. \textbf{3.} Il indexe le contenu pour la recherche vectorielle (RAG). \textbf{4.} Il publie ou archive les articles. \textbf{5.} Les agents IA disposent de la nouvelle connaissance lors de la prochaine requête.}
  {Conflit de version~: édition concurrente détectée~; fusion manuelle requise. Indexation échouée~: notification~; retry automatique toutes les 5 minutes.}
  {Articles publiés~; base de connaissances mise à jour~; agents IA bénéficiant de la nouvelle connaissance}

\ucfiche{UC-31}{Consulter les journaux système}
  {Administrator}
  {Accès Administrator~; journaux disponibles}
  {L'administrateur consulte les journaux système pour diagnostiquer des incidents ou auditer des actions}
  {\textbf{1.} L'administrateur accède au module Journaux. \textbf{2.} Il filtre par~: projet, utilisateur, type d'opération, plage de dates, niveau de log (INFO, WARNING, ERROR). \textbf{3.} Il consulte les entrées~: horodatage, acteur, action, payload, résultat. \textbf{4.} Il peut exporter les logs filtrés en CSV.}
  {Volume de logs élevé~: pagination et chargement progressif. Export > 10\,000 lignes~: génération asynchrone avec notification email.}
  {Journaux consultés~; export disponible~; aucune modification de l'état du système}

\ucfiche{UC-32}{Configurer les agents IA}
  {Administrator}
  {Accès Administrator~; clé API Claude disponible}
  {L'administrateur configure les paramètres des agents IA (modèle, contexte, permissions)}
  {\textbf{1.} L'administrateur accède au module Agents IA. \textbf{2.} Il configure les paramètres globaux~: modèle Claude (Opus/Sonnet), température, longueur max de contexte. \textbf{3.} Il définit les permissions par rôle~: quels agents sont accessibles à quels acteurs. \textbf{4.} Il active ou désactive les agents par projet. \textbf{5.} Il sauvegarde~; les changements prennent effet immédiatement.}
  {Clé API invalide~: message d'erreur~; test de connectivité disponible. Modèle non disponible~: fallback vers le modèle par défaut avec notification.}
  {Configuration des agents sauvegardée~; agents disponibles pour les acteurs autorisés dans leurs projets}

% =====================================================================
\section{Annexe B~--- Structure complète du bundle Databricks}
\label{annexe:bundle}
% =====================================================================

Cette annexe détaille la structure complète d'un bundle Databricks généré par ADMP~AI~Studio, à titre d'exemple pour un projet de migration standard composé de trois unités de migration (\texttt{CUS01}, \texttt{ACC00}, \texttt{APD}).

\subsection*{Arborescence du bundle}

\begin{tcolorbox}[colback=accenturegray!40, colframe=accenturedark, boxrule=0.5pt, arc=3pt,
  left=8pt, right=8pt, top=5pt, bottom=5pt, fontupper=\small\ttfamily]
poc\_\textit{project}\_YYYYMMDD\_bundle/\newline
\quad poc\_\textit{project}\_YYYYMMDD/\newline
\quad\quad runner.py\newline
\quad\quad validate.py\newline
\quad\quad MANIFEST.json\newline
\quad\quad units/\newline
\quad\quad\quad mig\_cus01.py\newline
\quad\quad\quad mig\_acc00.py\newline
\quad\quad\quad mig\_apd.py\newline
\quad\quad framework/\newline
\quad\quad\quad unit\_base.py\newline
\quad\quad config/\newline
\quad\quad\quad parameters.py\newline
\quad\quad\quad transco.py\newline
\quad\quad source\_ddl/\newline
\quad\quad\quad create\_tables.py\newline
\quad\quad\quad generate.py\newline
\quad\quad reports/\newline
\quad\quad\quad generate\_reports.py\newline
\quad\quad\quad integrity\_checks.json\newline
poc\_\textit{project}\_YYYYMMDD.whl
\end{tcolorbox}

\subsection*{Rôle de chaque composant}

\begin{table}[H]
\centering
\caption{Description des composants du bundle Databricks}
\label{tab:bundle-detail}
\begin{tabularx}{\textwidth}{llX}
\toprule
\textbf{Fichier} & \textbf{Type} & \textbf{Rôle} \\
\midrule
\texttt{runner.py} & Orchestrateur &
  Point d'entrée principal~; exécute les unités dans l'ordre configuré~; passe les paramètres run-time (catalog, schéma, run\_id) \\
\addlinespace
\texttt{validate.py} & Validation pré-run &
  Vérifie l'existence des tables cibles dans Unity Catalog~; contrôle la cohérence des paramètres avant toute écriture \\
\addlinespace
\texttt{MANIFEST.json} & Méta-données &
  Contient la version du bundle, les unités incluses, la date de génération et la version du compilateur \\
\addlinespace
\texttt{units/mig\_*.py} & Unités PySpark &
  Chaque fichier contient la logique de migration d'une entité~: DataFrame Spark, UNION ALL des branches, LEFT JOIN transco, INSERT INTO Delta \\
\addlinespace
\texttt{framework/unit\_base.py} & Classe de base &
  Fournit les mécanismes de logging (compteurs lignes lues/écrites/rejetées), gestion des erreurs et accès aux paramètres \\
\addlinespace
\texttt{config/parameters.py} & Paramètres &
  Dictionnaire typé des paramètres métier et des références catalog/schéma, modifiable avant chaque run \\
\addlinespace
\texttt{config/transco.py} & Transcodification &
  Charge les tables de transco depuis les fichiers de référence et les diffuse en broadcast Spark pour optimiser les jointures \\
\addlinespace
\texttt{source\_ddl/create\_tables.py} & DDL source &
  Recrée les tables sources en mémoire (tables Delta temporaires) pour permettre une exécution offline sans connexion aux systèmes sources \\
\addlinespace
\texttt{source\_ddl/generate.py} & Données de test &
  Génère des données de test synthétiques pour les tables sources~; utilisé dans les environnements de validation (UC-20) \\
\addlinespace
\texttt{reports/generate\_reports.py} & Rapports Excel &
  Génère les rapports de contrôle d'intégrité référentielle par unité et le rapport consolidé global \\
\addlinespace
\texttt{reports/integrity\_checks.json} & Définitions &
  Contient la définition des contrôles d'intégrité à exécuter~: clés primaires, clés étrangères, unicité, règles de gestion \\
\bottomrule
\end{tabularx}
\end{table}

% =====================================================================
\section{Annexe C~--- Modèle de données PostgreSQL}
\label{annexe:schema}
% =====================================================================

Cette annexe présente les principales tables de la base de données PostgreSQL d'ADMP~AI~Studio, telles que définies dans les migrations Flyway du dépôt. Les colonnes de verrouillage de section (checkout) sont détaillées car elles constituent un mécanisme architectural clé.

\begin{table}[H]
\centering
\caption{Table \texttt{migration\_units}~--- colonnes de verrouillage de section}
\label{tab:migration-units-lock}
\begin{tabularx}{\textwidth}{llX}
\toprule
\textbf{Colonne} & \textbf{Type} & \textbf{Description} \\
\midrule
\texttt{id} & UUID & Identifiant unique de l'unité de migration \\
\texttt{project\_id} & UUID & Référence au projet parent (FK) \\
\texttt{name} & VARCHAR & Nom de l'unité (ex.~\texttt{CUS01}, \texttt{ACC00}) \\
\texttt{data\_structure\_locked\_by} & UUID & ID utilisateur détenant le verrou sur la section Data Structure~; \texttt{NULL} si section libre \\
\texttt{record\_creation\_locked\_by} & UUID & ID utilisateur détenant le verrou sur la section RC \\
\texttt{data\_population\_locked\_by} & UUID & ID utilisateur détenant le verrou sur la section DP \\
\texttt{checkout\_baseline} & INTEGER & Numéro de version de référence au moment du checkout~; utilisé pour détecter les conflits \\
\texttt{created\_at} & TIMESTAMPTZ & Date de création de l'unité \\
\texttt{updated\_at} & TIMESTAMPTZ & Date de dernière modification \\
\bottomrule
\end{tabularx}
\end{table}

Le mécanisme de verrouillage garantit l'exclusivité d'édition par section~: chaque colonne \texttt{*\_locked\_by} peut être \texttt{NULL} (section libre) ou contenir l'UUID de l'utilisateur qui détient le verrou. La colonne \texttt{checkout\_baseline} permet de détecter les modifications concurrentes~: si la version courante de l'unité diffère de la baseline au moment du check-in, un conflit est signalé et une résolution manuelle est demandée.
